In this work, we presented DISCO-TAB, a unified framework for biomedical tabular data synthesis that bridges the gap between the interpretability of rule-based systems and the generative power of LLMs. By orchestrating a fine-tuned GPT model with a hierarchical, multi-objective discriminator system, we demonstrated that it is possible to synthesize clinical records that are not only statistically faithful but also structurally and semantically valid.

Our evaluation across \rev{nine} benchmarks \rev{indicates that DISCO-TAB is competitive with, and on several datasets improves upon, strong} GAN, \rev{diffusion,} and \rev{agentic} baselines, particularly in the challenging regimes of small sample sizes and severe class imbalance common to medical research. The introduction of Inverse-Frequency Reward Shaping \rev{is designed to keep} rare pathological cases \rev{from being ignored}, while our Automated Constraint Discovery preserves \rev{the pairwise clinical dependencies} essential for trustworthy decision support. \rev{We emphasize that DISCO-TAB provides no formal differential-privacy guarantee: the privacy-related results reported here are empirical, proximity-based disclosure-risk indicators rather than mathematically bounded protections, and deployment on sensitive data should be accompanied by an appropriate disclosure-risk assessment.}

As healthcare moves towards data-centric AI, DISCO-TAB offers a \rev{practical} mechanism to \rev{help} unlock the value of siloed EHR repositories. \rev{Future work will focus on three concrete directions: (i) integrating a formal differential-privacy mechanism (e.g., DP-SGD or a PATE-style aggregation) into the RL reward loop to provide bounded privacy guarantees; (ii) extending Automated Constraint Discovery beyond pairwise statistics toward higher-order and causal dependencies, ideally with clinician-in-the-loop validation of clinical plausibility; and (iii) validating the framework on large, multi-institutional EHR cohorts and extending the hierarchical discriminator to longitudinal, multi-modal patient trajectories.}